\chapter{Getting past obstacles}
\label{ch:t-obstacles}

A substation's walkways are narrow and lined with energized equipment: getting past a bucket or a
person must never mean walking off the walkway. The X30 has five layers, from the gentlest to the
last resort:

\begin{center}
\begin{tabularx}{\linewidth}{@{}>{\bfseries}p{0.22\linewidth}X@{}}
  \toprule
  Layer & What it does \\
  \midrule
  Path shift & slides sideways off the route and back, keeping the heading (holonomic robots) \\
  Detour & stuck: turns toward a gap and walks a short straight through it \\
  Backing up & stuck for long: backs up along the driven path and tries again \\
  The guard & the follower stops its own command before it could hit anything \\
  The field & the control mux stops whatever is left (Chapter~\ref{ch:t-control}) \\
  \bottomrule
\end{tabularx}
\end{center}

\section{Path shift}

The free corridor along the route ahead (\param{look_ahead}, \SI{2}{m} beyond the front edge) is
evaluated in route coordinates for candidate lateral offsets every \param{step} (\SI{0.1}{m}) up to
\param{max_offset} (\SI{0.8}{m}) on either side, the body plus \param{clearance} beside it. The
robot takes the free offset nearest the one it holds: a new side must be clearer by
\param{switch_margin} and pays \param{switch_cost} (hysteresis), so it does not hesitate between
left and right. It slides onto the offset at \param{offset_rate}, and comes back onto the route once
the way back has stayed free for \param{return_delay_sec}.

When no offset is free, the robot slows to the speed that stops \param{stop_margin} before the
obstacle (Eq.~\ref{eq:t-stopping-speed} on the free distance less the margin) and waits. After
\param{patience_sec} it searches wider, up to \param{search_offset} (\SI{1.5}{m}). A lane with less
than \param{centre_room} beside the body is walked in its middle.

\begin{dangerbox}[Set the offsets to the walkway]
\param{max_offset} and \param{search_offset} are how far the robot may leave the taught route.
Next to energized equipment they must be smaller than the room the walkway gives on either side,
less the minimum approach distance. A walkway \SI{1.5}{m} wide taught along its middle leaves
$(1.5 - 0.47)/2 \approx \SI{0.5}{m}$: set \param{max_offset} and \param{search_offset} to
\SI{0.5}{m} or less there, or teach the route where there is room.
\end{dangerbox}

\section{Detours and backing up}

\begin{description}
  \item[Detour] stuck for \param{detour.after_sec} (\SI{1.5}{s}) while following, the robot turns
    toward a gap within \param{max_angle_deg} (\ang{75}) of the route's direction and walks straight
    through it, at most \param{length} (\SI{2}{m}) at \param{speed}, only if the run brings it
    \param{min_progress} along the route. A detour leaves the route more than path shift does:
    disable it (\cmd{detour.enabled: false}) where the walkway leaves no room.
  \item[Backing up] no progress along the route for \param{backtrack.after_sec} (\SI{30}{s}): the
    robot backs up \param{distance} (\SI{1}{m}) along the path it walked, with the guard on the
    rear sectors, and tries again; after \param{max_attempts} at one place it stays BLOCKED and
    NETRA gets the alert.
\end{description}

\section{Reactive avoidance}

\code{obstacle_avoidance_node} reads the front sector scan and turns it into a bend of the
command: a heading offset, a sideways bias and a speed scale. With \cmd{avoidance_mode: reactive}
the follower applies them; in every mode but \code{off}, its \textbf{emergency} stops the forward
motion: an obstacle within \param{emergency_distance} (\SI{0.75}{m}) in the frontal
\param{emergency_cone_deg} (\ang{60}). With path shift and \param{emergency_slide}, the robot may
still slide out sideways.

Every policy works on the danger of each sector, from 0 beyond \param{safe_distance} $r_s$
(\SI{1.5}{m}) to 1 at the robot:
\begin{equation}
  \delta_j = \operatorname{clip}\!\left(1 - \frac{r_j}{r_s},\ 0,\ 1\right)
\end{equation}

\begin{description}
  \item[\code{vfh}] (the X30's) the headings whose sectors are free, in runs at least
    \param{min_gap_width_deg} wide; it takes the free heading nearest straight ahead and slows by
    the frontal danger: $s = \max(0, 1 - k_v\,\delta_{\text{front}})$.
  \item[\code{follow_gap}] removes a bubble of \param{bubble_radius} around the nearest obstacle
    (half-angle $\arctan(r_b/r_{\min})$) and aims at the centre of the widest remaining gap.
  \item[\code{potential_field}] each obstacle inside $r_s$ pushes with
    $F = k_r\,(1/r - 1/r_s)/r^2$; the sideways sum steers ($k_\theta\arctan F_y$) and pushes
    (\param{lateral_gain}, capped at \param{max_lateral}).
  \item[\code{braitenberg}] the danger left and right (triangular weights around
    \param{steering_zone_center_deg}) steers away: $\Delta\theta = k_s(\delta_R - \delta_L)$; the
    danger ahead and beside slows: $s = \max(0, 1 - k_v\delta_{\text{side}})(1 - \delta_{\text{ahead}})$.
\end{description}

The outputs are smoothed with \param{shared.smoothing} toward steering and speeding up, never
toward braking: a stop is never delayed.

\section{The guard}

Before it publishes, the follower checks its own command against the safety scan with the same
braking and separation model as the field, with its own margins (\param{safety.*},
Chapter~\ref{ch:t-following}). The guard's margins are at least the field's, and the follower
checks that at start: the follower stops itself before the field ever has to. A field that limits
the follower's commands in normal operation means the two are tuned apart.
